\section{Related work}
\label{sec:related}



\paragraph{Generalist robot manipulation policies.}
In recent years, robotic manipulation research has shifted focus from task-specific designs toward developing generalist policies capable of operating across diverse scenarios and embodiments~\citep{berscheid2019robot,brohan2022rt,dasari2019robonet,ebert2021bridge,fang2023rh20t,jang2022bc,khazatsky2024droid,mandlekar2018roboturk,walke2023bridgedata,shafiullah2023bringing}.
A representative approach is the development of vision-language-action (VLA) models.
In this paradigm, vision-language models (VLMs)---pretrained on vast, web-scale multimodal data (e.g., image, video, text)---are further trained on collection of diverse robotics datasets~\citep{black2410pi0, brohan2022rt,zheng2024tracevla,zhao2025cot,team2025gemini,li2024cogact,qu2025spatialvla,kim2024openvla}.
This knowledge transfer significantly reduces the reliance on domain-specific robotics data, which is typically expensive to collect via teleoperation.
However, with the notable exception of our previous work, MolmoAct~\citep{lee2025molmoactactionreasoningmodels}, reproducing frontier VLA models (e.g., the $\pi$ series~\citep{black2410pi0,intelligence2025pi05visionlanguageactionmodelopenworld})
remains highly challenging due to closed-source training data and code.
\molmoact{} represents our latest effort to democratize VLA research by releasing full training code, data, and model checkpoints that achieve state-of-the-art performance across a wide range of robotic manipulation benchmarks.

\paragraph{Embodied reasoning for robotic manipulation.}
While textual chain-of-thought (CoT) prompting~\citep{wei2022chain} has advanced reasoning in domains like program synthesis and mathematics, its benefits for robotic manipulation remain limited, likely because tacit physical knowledge is rarely articulated in the text corpora used to train backbone LLMs.
This has motivated a growing body of work on non-textual `embodied CoT' that leverages visual or spatial representations, including object bounding boxes and points~\citep{zawalski2024robotic,yuan2024robopoint,li2025hamster}, future image predictions~\citep{zhao2025cot}, latent reasoning~\citep{tur2026recurrent}, and 2D/3D point trajectories~\citep{sun2024emma,huang2025thinkact,lee2025molmoactactionreasoningmodels}.
Closest to our adaptive depth approach is PEEK~\citep{zhang2025peek}, which instead masks the RGB observation and conditions on visual tracing to learn a low-level policy.
A common drawback of these methods is heavy token consumption before action generation, which severely degrades inference latency.
\molmothink addresses this with adaptive depth reasoning: new depth tokens are computed only for novel visual signals, substantially improving latency, especially in third-person viewpoint setups.

\paragraph{Bridging VLM and action expert.}
While the architecture of VLA models has evolved significantly, the interface between VLM and action generators remains an open research problem.
Recent approaches either fully discretize actions via vector quantization~\citep{brohan2022rt,kim2024openvla,pertsch2025fast} or append a continuous generative action expert with diffusion~\citep{chi2025diffusion} or flow-matching objective~\citep{Lipman2022FlowModeling,black2410pi0}.
Standard continuous action experts typically condition only on the VLM's final hidden states.
\molmoactt{} departs from this by introducing a deep, layer-wise conditioning interface.
After pre-training the VLM on discrete action tokens produced by \molmoactfast for seamless multimodal data mixing, we post-train a continuous flow-matching expert using per-layer KV conditioning.
Unlike final-layer conditioning ~\citep{gr00tn1_2025,wang2026vla}, this per-layer KV conditioning gives the action expert dense access to the backbone's hierarchical visual-semantic features.
Finally, to explicitly ground this reasoning in 3D space without bottlenecking latency, \molmoactt-Think precedes this continuous action generation with an adaptive generation of discrete depth tokens only for dynamic scene regions.
